\documentclass[10pt,a4paper]{amsart}
\usepackage[margin=19mm]{geometry}
\usepackage{amsmath,amssymb,mathtools}
\usepackage[scr=rsfs,cal=euler]{mathalfa}
\usepackage{xfrac}
\usepackage[unicode,hidelinks,bookmarks=false]{hyperref}
\newcommand{\ie}{i.e.\ }
\newcommand{\cf}{cf.\ }
\newcommand{\resp}{respectively\ }
\newcommand{\Subsection}{\subsection}
\providecommand{\nobreakdash}{\nobreak\hbox{-}}
\setlength{\emergencystretch}{2em}
\multlinegap0pt
\usepackage{bm}
\usepackage{enumerate}
\usepackage{amssymb}
\usepackage[shortlabels]{enumitem}
\usepackage{xcolor}
\usepackage{dsfont}
\usepackage{mathtools}
\usepackage[normalem]{ulem}
\usepackage{tikz}
\newtheorem{df}{Definition}
\newtheorem{thm}{Theorem}
\newcommand{\modar}{\color{black}}
\title{On the role of symmetry for staircase mechanisms in local differential privacy efficiency across different privacy regimes}
\author{Chiara Amorino, Arnaud Gloter}
\date{Source: arXiv:2603.27572v1 (CC BY 4.0)}
\begin{document}
\maketitle

\noindent\emph{Source and licence.} On the role of symmetry for staircase mechanisms in local differential privacy efficiency across different privacy regimes. Version arXiv:2603.27572v1 (CC BY 4.0), distributed by arXiv under the Creative Commons Attribution 4.0 International licence, http://creativecommons.org/licenses/by/4.0/, as recorded in the arXiv licence metadata for this exact version and shown on the abstract page https://arxiv.org/abs/2603.27572v1. Full licence notice: \texttt{licenses/arxiv-2603.27572-CC-BY-4.0.txt}.\par\smallskip

\noindent\textit{Editorial scope.} Counted unit: the article's thm sym (source0.tex lines 801--810). The article's complete original proof of it is delivered with it (source0.tex lines 1661--1723, one \texttt{proof} environment of 63 lines, which the article places away from the statement). No mathematics is written, repaired or re-derived: the statement and the proof are the article's own lines, in the article's own notation, and no retained line is substituted or re-flowed.  Retained uncounted as the source-local context the proof needs: lines 544--546; lines 561--563; lines 657--669; lines 671--677; lines 701--745; lines 1410--1410; lines 1511--1513; lines 1521--1524; lines 1534--1537.  Nothing was omitted from any retained span, so the package carries no omission record.  No retained line contains a citation, so the wrapper carries no bibliography and no bibliography entry.  No retained line contains a figure environment, an image-inclusion command or drawing code, so no image is delivered and the package declares no asset dependency.  Editorial formatting only, in addition: an AMS \texttt{amsart} preamble that restates the article's own declarations and macros used by the retained lines, namely the environments \texttt{df}, \texttt{thm} on separate counters, together with the standard packages those lines need.  The retained environments print in this document as Df 1, Theorem 1 in order of appearance, whereas the article prints the same items with the numbering of its own full document.  The complete native source of the article as distributed in the arXiv e-print for this exact version is bundled unchanged as \texttt{sources/197358/source0.tex}.
\begin{equation} \label{eq : cond norm mu continue}
    \text{for a.e.\ } x \in \mathcal{X}, \quad \int_{\mathcal{E}} e_x(r) \, \mu(dr) = 1,
\end{equation}

\begin{equation}\label{eq: bound mu E}
    e^{-\alpha} \le \mu(\mathcal{E}) \le 1,
\end{equation}

\begin{equation}\label{eq : def i r ctn}
    i(r) := \frac{
        \left( \int_\mathcal{X} s_{\theta_0}(x) \, r(x) \, p_{\theta_0}(x) \, dx \right)^2
    }{
        \int_\mathcal{X} r(x) \, p_{\theta_0}(x) \, dx
    }
    = (e^\alpha - 1)^2
    \frac{
        \left( \int_{F_r^+} s_{\theta_0}(x) \, p_{\theta_0}(x) \, dx \right)^2
    }{
        1 + (e^\alpha - 1) \int_{F_r^+} p_{\theta_0}(x) \, dx
    },
\end{equation}

\begin{equation}\label{eq : optimisation pb ctn}
    I^* := \sup_{\mu}
    \int_\mathcal{E} i(r) \, \mu(dr),
    \quad
    \text{subject to } \int_\mathcal{E} e_x(r) \, \mu(dr) = 1
    \quad \text{for $dx$-almost every } x,
\end{equation}

To formalize this idea, let us consider Radon measures $\mu$ on $\mathcal{E}$ satisfying the normalization condition
\[
\int_\mathcal{E} e_x(r)\, \mu(dr) = 1 \quad \text{for a.e.\ } dx,
\]
and let us introduce the operator $T: \mathcal{E} \to \mathcal{E}$ defined by
\begin{equation}{\label{eq: def T sym}}
 T(r) := e^\alpha + 1 - r,   
\end{equation}
which essentially swaps the values $1$ and $e^\alpha$: for each $x$, we have $T(r)(x) = 1$ if and only if $r(x) = e^\alpha$. This operator captures a natural form of symmetry with respect to the extremal values in $\mathcal{E}$. Given a Radon measure $\mu$ on $\mathcal{E}$, we define the pushforward measure $T(\mu)$ by
\[
T(\mu) := \mu \circ T,
\]
so that for any integrable function $g$ on $\mathcal{E}$,
\[
\int_\mathcal{E} g(r)\, T(\mu)(dr) = \int_\mathcal{E} g(T(r))\, \mu(dr).
\]
{We begin by introducing the formal notion of symmetry that serves as a cornerstone for this work.
\begin{df}\label{def: sym mu}
Let $\mu$ be a Radon measure on $\mathcal{E}$ satisfying the normalization condition \eqref{eq : cond norm mu continue}. We say that $\mu$ is symmetric if it is invariant under the transformation $T$, i.e.,
\[
\mu = T(\mu).
\]
\end{df}
Although symmetry is defined via $T$-invariance, it is important to observe that $T(\mu)$ does not, in general, satisfy the normalization condition. To see this, we compute:
\begin{align*}
\int_\mathcal{E} e_x(r)\, T(\mu)(dr) 
&= \int_\mathcal{E} T(e_x(r))\, \mu(dr) 
= \int_\mathcal{E} \big(e^\alpha + 1 - e_x(r)\big)\, \mu(dr) \\
&= (e^\alpha + 1)\mu(\mathcal{E}) - \int_\mathcal{E} e_x(r)\, \mu(dr) .
%= (e^\alpha + 1)\mu(\mathcal{E}) - 1,
\end{align*}
%where the last equality follows from the normalization condition on $\mu$. 
%This naturally leads us to consider a rescaled version of $T(\mu)$ that restores the normalization. Define the correction factor
%\begin{equation}{\label{eq: c mu}}
% c(\mu) := (e^\alpha + 1)\mu(\mathcal{E}) - \int_\mathcal{E} e_x(r)\, \mu(dr).    
%\end{equation}
%We can now define the normalized pushforward measure
%\[ T^1(\mu) := \frac{T(\mu)}{c(\mu)}.\]
%By construction, $T^1(\mu)$ satisfies the required normalization:
%\[ \int_\mathcal{E} e_x(r)\, T^1(\mu)(dr) = \frac{1}{c(\mu)} \int_\mathcal{E} e_x(r)\, T(\mu)(dr) = \frac{c(\mu)}{c(\mu)} = 1.\]
This discrepancy suggests that any arbitrary measure must be appropriately transformed to satisfy both symmetry and normalization simultaneously. As it turns out, any Radon measure $\mu$—irrespective of its initial properties—can be mapped to such a state via the following construction:
\begin{equation}\label{eq: mu(s)}
\mu^{(s)} := \frac{1}{(e^\alpha + 1)\mu(\mathcal{E})} \big( \mu + T(\mu) \big).
\end{equation}
The fact that $\mu^{(s)}$ is both normalized and symmetric (in the sense of Definition \ref{def: sym mu}) is established in the following proposition, the proof of which is provided in Section \ref{s: proof main}.}

\section{Proofs}{\label{s: proof main}}

\begin{equation}\label{eq: fish step 1 3}
  \mathcal{I}_{\theta_0}(q^{(\mu)} \circ \mathcal{P}) = \int_\mathcal{E} i(r) \, \frac{1}{2}[\mu + T(\mu)](dr) = \frac{1}{2} \int_\mathcal{E} [i(r) + i(T(r))] \mu(dr),
\end{equation}

\begin{align}{\label{eq: i(T(r))}}
i(T(r)) &= i(e^\alpha + 1 - r) = \frac{\left( \int_\mathcal{X} s_{\theta_0}(x) (e^\alpha + 1 - r(x)) p_{\theta_0}(x) \, dx \right)^2}{\int_\mathcal{X} (e^\alpha + 1 - r(x)) p_{\theta_0}(x) \, dx} \\
&= \frac{\left( \int_\mathcal{X} s_{\theta_0}(x) r(x) p_{\theta_0}(x) \, dx \right)^2}{e^\alpha + 1 - \int_\mathcal{X} r(x) p_{\theta_0}(x) \, dx}, \nonumber
\end{align}

\begin{align}{\label{eq: 40.9}}
\mathcal{I}_{\theta_0}(q^{(\mu)} \circ \mathcal{P}) &= \frac{1}{2} \int_\mathcal{E} \left( \int_\mathcal{X} s_{\theta_0}(x) r(x) p_{\theta_0}(x) \, dx \right)^2 \left( \frac{1}{\int_\mathcal{X} r(x) p_{\theta_0}(x) \, dx} + \frac{1}{e^\alpha + 1 - \int_\mathcal{X} r(x) p_{\theta_0}(x) \, dx} \right) \mu(dr) \nonumber \\
&= \frac{1}{2}(e^\alpha + 1) \int_\mathcal{E} \frac{\left( \int_\mathcal{X} s_{\theta_0}(x) r(x) p_{\theta_0}(x) \, dx \right)^2}{\left( \int_\mathcal{X} r(x) p_{\theta_0}(x) \, dx \right) \left( e^\alpha + 1 - \int_\mathcal{X} r(x) p_{\theta_0}(x) \, dx \right)} \mu(dr).
\end{align}

\begin{thm}\label{thm: Fisher sym}
Assume that the maximizer $\bar{\mu}$ of the optimization problem \eqref{eq : optimisation pb ctn} is symmetric in the sense of Definition~\ref{def: sym mu}. Then, the function $r \mapsto i^{(s)}(r)$ is $\bar{\mu}$-almost everywhere constant, and the maximum Fisher information is given by
\begin{align*}
\mathcal{J}_{\theta_0}^{\max, \alpha} &= \frac{2}{e^\alpha + 1} (\max_{r \in B} i^{(s)}(r)), 
\end{align*}
with
\begin{equation}{\label{eq: i(s)r}}
i^{(s)}(r) = \frac{(e^\alpha - 1)^2}{2} \frac{(\int_{F_r^+}s_{\theta_0}(x)p_{\theta_0}(x) dx)^2}{(e^\alpha  -(e^\alpha - 1) \int_{F_r^+}p_{\theta_0}(x)dx)(1 + (e^\alpha - 1) \int_{F_r^+}p_{\theta_0}(x)dx)}.
\end{equation}
\end{thm}

\begin{proof}
We start by decomposing the measure \( \mu \) into its symmetric and asymmetric parts:
\begin{equation}\label{eq: start decomp 1}
\mathcal{I}_{\theta_0}(q^{(\mu)} \circ \mathcal{P}) = \int_{\mathcal{E}} i(r) \mu(dr) = \int_{\mathcal{E}} i(r) \mu^{(s)}(dr) + \int_{\mathcal{E}} i(r) \mu^{(as)}(dr).
\end{equation}
Since \( \mu^{(s)} \) is symmetric, the first term equals \( \int_{\mathcal{E}} i^{(s)}(r) \mu(dr) \), as shown in the proof of Theorem~\ref{thm: Fisher sym} (cf.~Equation~\eqref{eq: fish step 1 3}). For the asymmetric part, using the definition of \( \mu^{(s)} \) in \eqref{eq: mu(s)}, we write
\[
\mu^{(as)} = \mu - \mu^{(s)} = \frac{1}{(1 + e^\alpha) \mu(\mathcal{E})} \left( ((1 + e^\alpha) \mu(\mathcal{E}) - 1)\mu - T(\mu) \right).
\]
Substituting into \eqref{eq: start decomp 1} and using the definition of pushforward measure, we get
\begin{equation}\label{eq: asym decomp 2}
\int_{\mathcal{E}} i(r) \mu^{(as)}(dr) = \frac{1}{(1 + e^\alpha) \mu(\mathcal{E})} \int_{\mathcal{E}} \left[i(r)((1 + e^\alpha) \mu(\mathcal{E}) - 1) - i(T(r)) \right] \mu(dr).
\end{equation}
Using the definitions of \( i(r) \) and \( i(T(r)) \) in \eqref{eq : def i r ctn} and \eqref{eq: i(T(r))}, we compute:
\begin{align*}
& i(r)((1 + e^\alpha) \mu(\mathcal{E}) - 1) - i(T(r)) \\
&= \left( \int_\mathcal{X} s_{\theta_0}(x) r(x) p_{\theta_0}(x) dx \right)^2 \left( \frac{(1 + e^\alpha) \mu(\mathcal{E}) - 1}{\int_\mathcal{X} r(x) p_{\theta_0}(x) dx} - \frac{1}{e^\alpha + 1 - \int_\mathcal{X} r(x) p_{\theta_0}(x) dx} \right) \\
&= \frac{\left( \int_\mathcal{X} s_{\theta_0}(x) r(x) p_{\theta_0}(x) dx \right)^2 (1 + e^\alpha)\mu(\mathcal{E})}{\left( \int_\mathcal{X} r(x) p_{\theta_0}(x) dx \right)(e^\alpha + 1 - \int_\mathcal{X} r(x) p_{\theta_0}(x) dx)} \left( e^\alpha + 1 - \frac{1}{\mu(\mathcal{E})} - \int_\mathcal{X} r(x) p_{\theta_0}(x) dx \right).
\end{align*}
Recall that, from \eqref{eq: 40.9},
\begin{align*}
i^{(s)}(r) = \frac{e^\alpha + 1}{2} \frac{\left( \int_\mathcal{X} s_{\theta_0}(x) r(x) p_{\theta_0}(x) \, dx \right)^2}{\left( \int_\mathcal{X} r(x) p_{\theta_0}(x) \, dx \right) \left( e^\alpha + 1 - \int_\mathcal{X} r(x) p_{\theta_0}(x) \, dx \right)}.
\end{align*}
Recognizing its expression in the equation above, this implies
\[
\int_{\mathcal{E}} i(r) \mu^{(as)}(dr) = 2 \mu(\mathcal{E}) \int_{\mathcal{E}} \left[ e^\alpha + 1 - \frac{1}{\mu(\mathcal{E})} - \int_{\mathcal{X}} r(x) p_{\theta_0}(x) dx \right] i^{(s)}(r) \mu(dr).
\]
Now, since \( r(x) = 1 + (e^\alpha - 1)\mathbf{1}_{F_r^+}(x) \), we find
\[
\int_\mathcal{X} r(x) p_{\theta_0}(x) dx = 1 + (e^\alpha - 1)\int_{F_r^+} p_{\theta_0}(x) dx,
\]
so that
\begin{align*}
e^\alpha + 1 - \frac{1}{\mu(\mathcal{E})} - \int_\mathcal{X} r(x) p_{\theta_0}(x) dx &= e^\alpha - \frac{1}{\mu(\mathcal{E})} - (e^\alpha - 1)\int_{F_r^+} p_{\theta_0}(x) dx \\
&= \left( \frac{e^\alpha + 1}{2} - \frac{1}{\mu(\mathcal{E})} \right) - (e^\alpha - 1)\left( \int_{F_r^+} p_{\theta_0}(x) dx - \frac{1}{2} \right).    
\end{align*}
It follows that
\begin{align*}
\mathcal{I}_{\theta_0}(q^{(\mu)} \circ \mathcal{P})& = \int_{\mathcal{E}} i^{(s)}(r) \mu(dr) + 2 \mu(\mathcal{E}) \left( \frac{e^\alpha + 1}{2} - \frac{1}{\mu(\mathcal{E})} \right) \int_{\mathcal{E}} i^{(s)}(r) \mu(dr) \\
& - 2 \mu(\mathcal{E})(e^\alpha - 1)\left( \int_{F_r^+} p_{\theta_0}(x) dx - \frac{1}{2} \right) \int_{\mathcal{E}} i^{(s)}(r) \mu(dr).
\end{align*}
Recalling that \( i^{(as)}(r) = i^{(s)}(r) \left( \int_{F_r^+} p_{\theta_0}(x) dx - \frac{1}{2} \right) \), this gives the desired decomposition.

Let us now prove that \( i^{(as)}(T(r)) = -i^{(as)}(r) \). Since \( i^{(s)}(T(r)) = i^{(s)}(r) \) and \( T \) exchanges \( F_r^+ \) with its complement, we have
\[
i^{(as)}(T(r)) = i^{(s)}(r) \left( \int_{(F_r^+)^c} p_{\theta_0}(x) dx - \frac{1}{2} \right) = i^{(s)}(r) \left(\frac{1}{2} - \int_{F_r^+} p_{\theta_0}(x) dx \right) = - i^{(as)}(r).
\]

To conclude, in the regime \( \alpha \to 0 \), Equation \eqref{eq: i(s)r} yields:
{\modar
	\begin{align*}
 i^{(s)}(r) &= \alpha^2 \left( \int_{F_r^+} s_{\theta_0}(x) p_{\theta_0}(x) dx \right)^2 + O(\alpha^3),
\\ 
\quad i^{(as)}(r) &=\alpha^2 \left( \int_{F_r^+} s_{\theta_0}(x) p_{\theta_0}(x) dx \right)^2 \left( \int_{F_r^+} p_{\theta_0}(x) dx - \frac{1}{2} \right)+O(\alpha^3).
\end{align*}}
Moreover, from the bound on \( \mu(\mathcal{E}) \) in \eqref{eq: bound mu E}, we find
\[{\modar
\frac{e^\alpha + 1}{2} - \frac{1}{\mu(\mathcal{E})} 
=\frac{e^\alpha + 1}{2} - \frac{1}{1+O(\alpha)}=O(\alpha), }
%\sim \frac{2 + \alpha}{2} - \frac{1}{1 - \alpha} = -\frac{\alpha}{2},
\]
which, combined with the decomposition above, yields the result.
\end{proof}

\end{document}
